\section{Exponential path integrability under the exact microscopic event}
\label{sec:v53-exponential-paths}
The fourth-moment argument can be strengthened, but only if its
constants are kept explicit in the differentiation order. We do so on
fixed complex neighborhoods of the collision spectrum. This yields an
exponential moment of the supremum of the entire pinned path, not an
exponential moment of the roof density and not a high-frequency estimate.

\subsection{Cauchy bounds of every order}
\begin{lemma}[Factorial bounds at the moving peaks]
\label{lem:v53-factorial-derivatives}
On a fixed spectral band, reduce the finitely many peak charts once.
There are $C,D,a>0$ and $\rho<1$, independent of the derivative order
$q\ge0$, the collision block length $L\ge1$, the radius and the peak,
such that, for every real unit vector $e$,
\begin{align}
 &\left\|\left.\frac{d^q}{dz^q}
 \left(e^{-iz e\cdot\mu_j(R)}\mathscr T_{R,\theta+ze}\right)^L
                                   \right|_{z=0}\right\|\notag\\
 &\qquad\le Cq!D^qL^{q/2}
       \left(e^{-aL|\theta-a_j(R)|^2}+\rho^L\right).
                                      \label{eq:v53-factorial-derivatives}
\end{align}
On the compact nonresonant complement the uncentered derivative is
bounded by $Cq!D^q\rho^L$. At $L=0$ the derivative is zero for
$q>0$ and the operator is the identity for $q=0$.
\end{lemma}
\begin{proof}
The proof of Lemma~\ref{lem:v41-peak-derivatives} provides more than
its stated fourth-order estimate. On the complex circle
$|z|=d_0/\sqrt L$ it bounds the centered principal power by
$C\exp(-aL|\theta-a_j|^2)$ and the complementary power by
$C\rho_0^L\exp(C\sqrt L)\le C'\rho^L$, with
$\rho_0<\rho<1$. The projection is bounded on the same circle.
The circle, the analytic neighborhood and these upper bounds do not
depend on $q$. Cauchy's formula therefore gives exactly the factor
$q!(\sqrt L/d_0)^q$. Taking $D\ge d_0^{-1}$ proves
\eqref{eq:v53-factorial-derivatives}. A fixed complex circle on the
nonresonant complement gives its asserted bound in the same way.
The finitely many charts and local parameter spaces permit common
constants. No claim is made about their growth with the fixed band.
\end{proof}

Fix a bounded roof interval $J$ of positive length and the original
event $E=E_{n,m,R}(k,t;J)$ in \eqref{eq:v41-exact-conditioning}.
The vector $\mathsf S_j$ keeps the displacement, roof and occupation
order of Section~\ref{sec:v41-bridges}.

\begin{proposition}[All pinned increment moments with factorial constants]
\label{prop:v53-all-increment-moments}
There is $C_J<\infty$, independent of the integer $r\ge1$, such that
for $0\le a<a+l\le m$,
\begin{equation}\label{eq:v53-all-increment-moments}
 \int_E\left|\mathsf S_{a+l,R}-\mathsf S_{a,R}
                         -\frac lm\mathsf S_{m,R}\right|^{2r}d\nu_R^*
            \le (2r)!C_J^{2r}l^r m^{-2}.
\end{equation}
The estimate is uniform over all original exact labels, roof targets
and radii. Neither $J$ nor an auxiliary spectral band depends on $m$
or on $r$.
\end{proposition}
\begin{proof}
Majorize $\one_J$ by one fixed nonnegative band-limited function with
integrable Fourier transform, as in
Proposition~\ref{prop:v41-pinned-fourth}. This increases the positive
moment integrand. Put $\alpha=l/m$ and $b=m-a-l$. In one coordinate,
$(-1)^r$ times the derivative of order $2r$ at zero of
\[
 \ell\left(M_{\eta_R}
 \mathscr T_{R,\theta-\alpha ze}^{b}
 \mathscr T_{R,\theta+(1-\alpha)ze}^{l}
 \mathscr T_{R,\theta-\alpha ze}^{a}(\eta_R\nu)\right)
\]
is the required Fourier-transformed moment. Differentiation under the
finite-time integral is justified for every fixed order by boundedness
of the single-collision record. We bound its absolute value.

On each moving peak, the centering factors cancel identically since
$-\alpha(a+b)+(1-\alpha)l=0$. By
Lemma~\ref{lem:v53-factorial-derivatives}, Leibniz's formula has total
size at most
\[
 C(2r)!D^{2r}
       \{\alpha(\sqrt a+\sqrt b)+(1-\alpha)\sqrt l\}^{2r}
       \left(e^{-a_1m|\theta-a_j|^2}+\rho^{m/3}\right).
\]
To see the factorial dependence, the multinomial coefficient cancels
the three derivative factorials, leaving $(2r)!$ times the sum of
monomials in the three displayed size factors. That sum is at most
their sum raised to $2r$. In every spectral term at least one block
has length at least $m/3$; its Gaussian or complementary factor gives
the final parenthesis. The other two blocks have bounded decay
factors, including the zero-length convention. Since
$\alpha(\sqrt a+\sqrt b)\le\sqrt2\,l/\sqrt m\le\sqrt{2l}$,
the size is at most $(2r)!C^{2r}l^r$ times that parenthesis.
The Gaussian integral is $O(m^{-2})$ in the four collision frequencies.

On the compact nonresonant part, use a single fixed complex disk for
the entire three-block product. All its powers remain exponentially
contracting there; its norm is at most $C\rho^m$. Cauchy's formula
gives $C(2r)!D^{2r}\rho^m$, bounded by
$(2r)!C_J^{2r}l^rm^{-2}$ for all $l\ge1$. Integrate against the
fixed roof majorant's Fourier transform and sum the finite charts.
Passing from coordinates to the Euclidean norm costs only a factor
of the form $C^{2r}$. The original source factor $c^{-1}$ is fixed
and is absorbed in $C_J$, not reintroduced as a changed probability.
\end{proof}

\subsection{A maximal estimate before conditioning}
For $I_t=[t,t+h]$, $h>0$, write
\[
 \lambda_{m,R}^{n,k,t}=(m^2/h)\one_E\nu_R^*,\qquad
 E=\{K_{n,R}=k,N_{n,R}=m,T_{n,R}\in I_t\}.
\]
Its total mass is bounded uniformly, by the fixed-window local upper
bound, even on label sequences for which the mass tends to zero.

\begin{theorem}[Exponential supremum moments of the complete paths]
\label{thm:v53-exponential-maximal}
For every fixed $h>0$ there are $a_h,C_h>0$ such that
\begin{equation}\label{eq:v53-collision-exponential}
 \sup_{m\ge2,R,n,k,t}\int
     \exp(a_h\|\mathcal B_{m,R}\|_\infty)\,d\lambda_{m,R}^{n,k,t}
                                                           \le C_h.
\end{equation}
For every fixed $M_0<\infty$ there are $a_{h,M_0},C_{h,M_0}>0$
such that, on $|Z(t)|\le M_0$ and for all sufficiently large $m$,
\begin{equation}\label{eq:v53-return-exponential}
 \int\exp\!\left(a_{h,M_0}
       (\|\mathcal B_{m,R}\|_\infty+\|\mathcal Y_{n,R}\|_\infty)\right)
                                      d\lambda_{m,R}^{n,k,t}\le C_{h,M_0}.
\end{equation}
No physical or section guard is imposed in these integrals.
\end{theorem}
\begin{proof}
Normalize \eqref{eq:v53-all-increment-moments} by the path scale and
by $m^2/h$. Interpolation gives, for all real $s,t\in[0,1]$ and
all even $q\ge2$,
\[
 \int|\mathcal B_{m,R}(t)-\mathcal B_{m,R}(s)|^q\,d\lambda
                         \le q!C_h^q|t-s|^{q/2}.
\]
Within one grid interval the fractional interpolation factor is of
order $|t-s|^q m^q$, which gives this bound. For a general increment
split off its two fractional end pieces and the intervening grid
interval; Minkowski's inequality changes $C_h$, independently of $q$.

For even $q\ge4$, the $L^q(\lambda)$ norm of the maximum of the
$2^j$ dyadic increments of length $2^{-j}$ is at most
\[
 (q!)^{1/q}C_h2^{-j(1/2-1/q)}.
\]
Every continuous path starting at zero is bounded in supremum by the
sum of these maxima over dyadic levels. Minkowski's inequality and
$1/2-1/q\ge1/4$ therefore yield
\[
 \int\|\mathcal B_{m,R}\|_\infty^q\,d\lambda\le q!C_h'^q,
                                                        \qquad q\ge4.
\]
The finite mass bound pays the smaller orders. Cauchy--Schwarz between
consecutive even moments pays odd orders; the ratio
$\sqrt{(q-1)!(q+1)!}/q!=\sqrt{(q+1)/q}$ is bounded.
After one enlargement the factorial bound holds at every integer
order. Expanding the exponential with $a_hC_h'<1$ proves
\eqref{eq:v53-collision-exponential}. This is an argument for finite
measures and does not divide by a rare-event probability.

On the central event, $n/m=c+O_{M_0}(m^{-1/2})$ and
$|\mathcal X_{m,R}(1)|\le M_0+h/\sqrt m$. The exact identities
\eqref{eq:v41-pinned-time-change} and
\eqref{eq:v52-exact-clock-displacement} give, at every return grid
point $l/n$,
\[
 |\mathcal Y_{n,R}(l/n)|
 \le\|\sqrt{m/n}L_R^{-1}\|
       \left(1+\frac{\sqrt m}{n}|\mathcal X_{m,R}(1)|\right)
                                     \|\mathcal B_{m,R}\|_\infty.
\]
The coefficient is uniformly bounded for large $m$. Linear
interpolation on the return grid cannot enlarge the maximum.
Reducing $a_{h,M_0}$ in \eqref{eq:v53-collision-exponential}
proves \eqref{eq:v53-return-exponential}. Thus the time change
preserves a quantitative moment budget on the unchanged source.
\end{proof}

\begin{corollary}[Exponentially weighted physical-defect mass]
\label{cor:v53-exponential-defect}
On fixed central windows there are $a,C>0$ such that, for every fixed
sufficiently small $\varepsilon>0$,
\begin{align}
 &\limsup_m\sup_{R,n,k,t:\ |Z(t)|\le M_0}\frac{m^2}{h}
 \int_E(1-H_{m,R}^{\varepsilon})
 e^{a(\|\mathcal B_{m,R}\|_\infty+\|\mathcal Y_{n,R}\|_\infty)}d\nu_R^*
                                  \le C\varepsilon^{1/32}.
                                      \label{eq:v53-exponential-defect}
\end{align}
The same estimate holds for either disjoint first-physical-defect type.
It concerns the original positive incidence and clearance sources,
without transporting them across a physical boundary.
\end{corollary}
\begin{proof}
Use Cauchy--Schwarz in $\lambda$, noting that
$(1-H^{\varepsilon})^2\le1-H^{\varepsilon}$. The first factor is
bounded by $C\varepsilon^{1/32}$ after the collision limsup, by
\eqref{eq:v49-local-limsup}. Choose $2a$ below the exponential
constant in \eqref{eq:v53-return-exponential} to bound the second
factor. Each positive first-defect subsource is dominated by the
full physical remainder, proving its assertion. This remains an
integrated estimate; it supplies no essential-height estimate in $u$.
\end{proof}
